% ==============================================================================
% Report_Lab02.tex
% CSC15012 - Applications of Natural Language Processing in Industry
% Project 2: Fine-tuning Intent Detection Model with Banking Dataset
% Student: Le Dinh Minh Quan - 23127460
% ==============================================================================
\documentclass[12pt,a4paper]{article}

% ---- Packages ----
\usepackage[utf8]{inputenc}
\usepackage[T5]{fontenc}
\usepackage[vietnamese,english]{babel}
\usepackage{geometry}
\geometry{margin=2.5cm}
\usepackage{graphicx}
\usepackage{amsmath,amssymb}
\usepackage{booktabs}
\usepackage{enumitem}
\usepackage{hyperref}
\hypersetup{colorlinks=true, linkcolor=blue, citecolor=blue, urlcolor=blue}
\usepackage{xurl}  % allow URL line-breaking at almost any character
\usepackage{listings}

% Allow LaTeX to be more relaxed about line endings (avoid right-margin overflow)
\sloppy
\emergencystretch=4em
\usepackage{xcolor}
\usepackage{float}
\usepackage{caption}
\usepackage{subcaption}
\usepackage{fancyhdr}
\usepackage{titlesec}
\usepackage{array}
\usepackage{longtable}
\usepackage{multicol}

% ---- Code listing style ----
\lstdefinestyle{mystyle}{
    backgroundcolor=\color{gray!5},
    basicstyle=\ttfamily\small,
    breaklines=true,
    captionpos=b,
    commentstyle=\color{green!50!black},
    keywordstyle=\color{blue},
    stringstyle=\color{red!70!black},
    numbers=left,
    numberstyle=\tiny\color{gray},
    numbersep=8pt,
    frame=single,
    rulecolor=\color{gray!30},
    tabsize=4,
    showstringspaces=false,
}
\lstset{style=mystyle}

% ---- Header/Footer ----
\pagestyle{fancy}
\fancyhf{}
\fancyhead[L]{\small CSC15012 -- Applications of NLP in Industry}
\fancyhead[R]{\small Project 2}
\fancyfoot[C]{\thepage}

% ---- Convenient project links ----
% Long-form URL (auto line-break with xurl) – use inside body paragraphs
\newcommand{\repolink}{\url{https://github.com/ledinhminhquan/banking-intent-unsloth}}
\newcommand{\notebooklink}{\url{https://drive.google.com/file/d/1dd4BxqryUvpkOvmDkm8-ZWor0LHpVSZu/view?usp=sharing}}
\newcommand{\videolink}{\url{https://drive.google.com/file/d/1D9hAiMirLMdjMijR-znzDAdYw2LGtdqP/view?usp=sharing}}

% Short clickable label – use inside compact lists / tables to avoid margin overflow
\newcommand{\repolinkshort}{\href{https://github.com/ledinhminhquan/banking-intent-unsloth}{github.com/ledinhminhquan/banking-intent-unsloth}}
\newcommand{\notebooklinkshort}{\href{https://drive.google.com/file/d/1dd4BxqryUvpkOvmDkm8-ZWor0LHpVSZu/view?usp=sharing}{Open notebook on Google Drive}}
\newcommand{\videolinkshort}{\href{https://drive.google.com/file/d/1D9hAiMirLMdjMijR-znzDAdYw2LGtdqP/view?usp=sharing}{Watch demo video on Google Drive}}

% ==============================================================================
\begin{document}

% ---- Title page ----
\begin{titlepage}
    \centering

    % University logo - upload your file as logo_hcmus.png
    \includegraphics[width=3cm]{logo_hcmus.png}

    \vspace{0.5cm}
    {\large \textbf{VIETNAM NATIONAL UNIVERSITY HO CHI MINH CITY}}\\[0.2cm]
    {\large \textbf{UNIVERSITY OF SCIENCE}}\\[0.2cm]
    {\large \textbf{FACULTY OF INFORMATION TECHNOLOGY}}\\[1.5cm]

    \rule{\textwidth}{0.4pt}\\[0.4cm]
    {\Huge \textbf{Project 2}}\\[0.3cm]
    {\LARGE \textbf{Fine-Tuning Intent Detection Model\\with Banking Dataset}}\\[0.2cm]
    \rule{\textwidth}{0.4pt}\\[1.5cm]

    {\Large \textbf{Course:} CSC15012 -- Applications of Natural Language Processing in Industry}\\[0.5cm]
    {\Large \textbf{Lecturer:} Dr. Nguyễn Hồng Bửu Long}\\[1.5cm]

    \begin{tabular}{ll}
        \textbf{Student:} & Lê Đình Minh Quân \\
        \textbf{Student ID:} & 23127460 \\
    \end{tabular}

    \vfill
    {\large Ho Chi Minh City, April 2026}
\end{titlepage}

% ---- Table of Contents ----
\newpage
\tableofcontents
\newpage

% ==============================================================================
\section{Introduction}
% ==============================================================================

This report documents the implementation and evaluation of a banking intent
classification system built for Project 2 of the course
\textit{CSC15012 -- Applications of Natural Language Processing in Industry}.
The objective is to fine-tune a Large Language Model (LLM) using
\textbf{Unsloth} and \textbf{QLoRA} (Quantized Low-Rank Adaptation) on the
\textbf{BANKING77} dataset, so that the model can accurately predict the intent
behind a customer banking-related message.

The BANKING77 dataset, introduced by Casanueva et al.~\cite{casanueva2020},
contains 13{,}083 annotated customer service queries across 77 fine-grained
intent categories that cover common banking operations such as card activation,
PIN changes, cross-border transfers, balance inquiries, exchange rate questions,
and lost or stolen card scenarios.

\subsection{Project Objectives}

The five tasks specified by the project brief are:
\begin{enumerate}[noitemsep]
    \item Sample and construct a suitable subset of the BANKING77 dataset.
    \item Preprocess the data and split it into training and testing sets.
    \item Fine-tune a text classification model using \textbf{Unsloth} and
          clearly document all hyperparameters used.
    \item Evaluate the fine-tuned model on an independent test set.
    \item Implement a standalone inference module that loads the saved
          checkpoint and predicts the intent of an input message.
\end{enumerate}

\subsection{Project Deliverables and Links}

\begin{itemize}[noitemsep]
    \item \textbf{GitHub Repository (public):} \repolinkshort{}
    \item \textbf{Google Colab Notebook (Drive):} \notebooklinkshort{}
    \item \textbf{Demo Video (Drive, public):} \videolinkshort{}
\end{itemize}

\noindent The full URLs are also provided here for offline reference:
\begin{itemize}[noitemsep]
    \item \repolink{}
    \item \notebooklink{}
    \item \videolink{}
\end{itemize}

\subsection{Final Headline Result}

The fine-tuned LLaMA 3.2 3B model achieved \textbf{87.92\% test accuracy}
($677 / 770$) on the held-out test set, with a macro F1-score of
\textbf{0.8779} across all 77 intent classes, after only \textbf{248.2 seconds}
of training on a single NVIDIA H100 GPU.

% ==============================================================================
\section{Data Preparation and Processing}
\label{sec:data}
% ==============================================================================

\subsection{Dataset Overview}

The BANKING77 dataset is hosted on the Hugging Face Hub at
\texttt{PolyAI/banking77}. The full dataset contains:
\begin{itemize}[noitemsep]
    \item \textbf{Training set:} 10{,}003 samples
    \item \textbf{Test set:} 3{,}080 samples
    \item \textbf{Number of intents:} 77 fine-grained categories
\end{itemize}

Example intent labels include
\texttt{activate\_my\_card}, \texttt{card\_delivery\_estimate},
\texttt{change\_pin}, \texttt{cancel\_transfer},
\texttt{lost\_or\_stolen\_card}, \texttt{exchange\_rate}, and
\texttt{extra\_charge\_on\_statement}, among others.

\subsection{Sampling Strategy}

To make training tractable on a single GPU within a Google Colab Pro session,
we sampled a stratified subset of the BANKING77 dataset:
\begin{itemize}[noitemsep]
    \item \textbf{Training:} up to 40 samples per intent class
          $\rightarrow$ \textbf{3{,}075 samples} total
          (a few classes have fewer than 40 examples in the original split).
    \item \textbf{Test:} up to 10 samples per intent class
          $\rightarrow$ \textbf{770 samples} total (10 per class for all 77 classes).
\end{itemize}
Sampling used a fixed random seed (\texttt{seed=42}) for reproducibility.

\subsection{Loading the Dataset Robustly}

Recent releases of the \texttt{datasets} library (>=3.0) no longer execute
arbitrary dataset loading scripts. Because the BANKING77 repo originally
shipped with a Python loading script, a naive
\texttt{load\_dataset("PolyAI/banking77")} call now raises
\textit{``Dataset scripts are no longer supported.''} To handle this, the
preprocessing module implements a three-tier fallback:
\begin{enumerate}[noitemsep]
    \item Try the standard \texttt{load\_dataset} call (works on older versions).
    \item If that fails, load from the auto-generated parquet branch by passing
          \texttt{revision} = \texttt{"refs/convert/parquet"}.
    \item If both fail, fetch the parquet files directly from the Hugging Face
          Hub URLs.
\end{enumerate}
This made the pipeline robust against future Hugging Face library changes.

\subsection{Preprocessing Steps}

\begin{enumerate}[noitemsep]
    \item \textbf{Text normalization:} all input text was stripped of leading
          and trailing whitespace, lower-cased, and had any internal whitespace
          collapsed into single spaces.
    \item \textbf{Label mapping:} a bidirectional dictionary
          (\texttt{id2label} / \texttt{label2id}) was constructed and saved to
          \texttt{sample\_data/label\_map.json}.
    \item \textbf{Format conversion:} the sampled splits were exported as
          UTF-8 CSV files with three columns: \texttt{text}, \texttt{label}
          (integer ID), and \texttt{label\_text} (human-readable intent name).
    \item \textbf{Shuffling:} both train and test splits were randomly
          shuffled with the same seed for reproducibility.
\end{enumerate}

\subsection{Data Statistics}

\begin{table}[H]
    \centering
    \caption{Dataset statistics after sampling and preprocessing.}
    \label{tab:data-stats}
    \begin{tabular}{lcc}
        \toprule
        \textbf{Split} & \textbf{Samples} & \textbf{Avg. tokens / sample} \\
        \midrule
        Training       & 3{,}075           & $\sim$12                      \\
        Test           & 770               & $\sim$12                      \\
        \bottomrule
    \end{tabular}
\end{table}

\begin{figure}[H]
    \centering
    \includegraphics[width=0.95\textwidth]{figs/01_sample_train_csv.png}
    \caption{First ten rows of \texttt{sample\_data/train.csv} after preprocessing.
    Each row contains the cleaned customer message, its integer label, and the
    human-readable intent name.}
    \label{fig:sample-data}
\end{figure}

% ==============================================================================
\section{Fine-Tuning with Unsloth}
\label{sec:finetuning}
% ==============================================================================

\subsection{Model Selection}

We used the \textbf{LLaMA 3.2 3B Instruct} model loaded in 4-bit quantization
through Unsloth's pre-quantized release:
\begin{center}
    \texttt{unsloth/Llama-3.2-3B-Instruct-bnb-4bit}
\end{center}

This model was selected for the following reasons:
\begin{itemize}[noitemsep]
    \item \textbf{Strong instruction-following baseline:} the Instruct variant
          already understands prompt-style classification reasonably well.
    \item \textbf{Memory efficiency:} 4-bit (NF4) quantization reduces VRAM
          usage roughly $4\times$ compared with full precision, so even a
          small consumer GPU can fine-tune it.
    \item \textbf{Speed:} Unsloth's custom kernels make this model
          significantly faster to fine-tune than vanilla
          \texttt{transformers}~\cite{unsloth}.
\end{itemize}

\subsection{LoRA Configuration}

We used Low-Rank Adaptation (LoRA)~\cite{hu2021lora} so that only a small
fraction of parameters are updated. The LoRA adapters were attached to all
attention and feed-forward projection layers:

\begin{table}[H]
    \centering
    \caption{LoRA hyperparameters.}
    \label{tab:lora}
    \begin{tabular}{ll}
        \toprule
        \textbf{Parameter}     & \textbf{Value}                                        \\
        \midrule
        LoRA rank ($r$)        & 16                                                    \\
        LoRA alpha ($\alpha$)  & 16                                                    \\
        LoRA dropout           & 0.0                                                   \\
        Target modules         & \texttt{q\_proj, k\_proj, v\_proj, o\_proj}            \\
                               & \texttt{gate\_proj, up\_proj, down\_proj}              \\
        Bias                   & none                                                  \\
        Gradient checkpointing & Unsloth (memory-efficient)                            \\
        Random state           & 3407                                                  \\
        \bottomrule
    \end{tabular}
\end{table}

\subsection{Prompt Format}

We used \textbf{Supervised Fine-Tuning (SFT)} with an Alpaca-style prompt
template. Each training example was formatted as:

\begin{lstlisting}[language={},caption={Prompt template used both for training and inference (the trailing label is omitted at inference time).},label={lst:prompt}]
Below is an instruction that describes a task, paired with an input
that provides further context. Write a response that appropriately
completes the request.

### Instruction:
Classify the intent of the following banking customer message.
Reply with only the intent label, nothing else.

### Input:
{customer_message}

### Response:
{intent_label}<|end_of_text|>
\end{lstlisting}

At inference time, the trailing \texttt{\{intent\_label\}} part is removed,
the model then auto-regressively generates the predicted label.

\subsection{Training Hyperparameters}

\begin{table}[H]
    \centering
    \caption{Training hyperparameters used in this run.}
    \label{tab:training-hp}
    \begin{tabular}{ll}
        \toprule
        \textbf{Hyperparameter}         & \textbf{Value}                       \\
        \midrule
        Per-device batch size           & 8                                    \\
        Gradient accumulation steps     & 4                                    \\
        Effective global batch size     & 32                                   \\
        Number of epochs                & 3                                    \\
        Learning rate                   & $2 \times 10^{-4}$                   \\
        Optimizer                       & AdamW 8-bit                          \\
        LR scheduler                    & Linear decay with warmup             \\
        Warmup steps                    & 10                                   \\
        Weight decay                    & 0.01                                 \\
        Max sequence length             & 512                                  \\
        Mixed precision                 & bf16 (auto fp16 fallback)            \\
        Random seed                     & 42                                   \\
        Trainer                         & \texttt{trl.SFTTrainer} + \texttt{SFTConfig} \\
        \bottomrule
    \end{tabular}
\end{table}

\subsection{Training Platform}

Training was conducted on \textbf{Google Colab Pro} with an
\textbf{NVIDIA H100 80GB HBM3} GPU. Total wall-clock training time was
\textbf{248.2 seconds} ($\approx 4.1$ minutes) for 3 epochs, with a final
training loss of \textbf{0.5320}.

\begin{figure}[H]
    \centering
    \includegraphics[width=0.95\textwidth]{figs/02_training_log.png}
    \caption{Training console output showing epoch-by-epoch loss decrease,
    final training loss, and total wall-clock training time.}
    \label{fig:training-loss}
\end{figure}

% ==============================================================================
\section{Evaluation Results}
\label{sec:results}
% ==============================================================================

\subsection{Test-Set Performance}

After fine-tuning, the model was evaluated on the held-out test set
(770 samples, 10 per intent class).

\begin{table}[H]
    \centering
    \caption{Overall test set performance.}
    \label{tab:results}
    \begin{tabular}{lc}
        \toprule
        \textbf{Metric}      & \textbf{Value}                  \\
        \midrule
        Test accuracy        & \textbf{87.92\%} ($677 / 770$)  \\
        Macro precision      & 0.8915                          \\
        Macro recall         & 0.8792                          \\
        Macro F1-score       & 0.8779                          \\
        Weighted F1-score    & 0.8779                          \\
        Training loss        & 0.5320                          \\
        Training time        & 248.2s ($\approx$ 4.1 min)      \\
        Evaluation time      & 125.6s for 770 samples          \\
        \bottomrule
    \end{tabular}
\end{table}

This is a strong result given that:
\begin{itemize}[noitemsep]
    \item The model is generative, not a dedicated classification head, so it
          must produce the exact intent label as text.
    \item Only 40 training samples per class were used, instead of the full
          training set.
    \item The model has 77 fine-grained classes, many of which are subtle
          variants of one another (e.g., \texttt{lost\_or\_stolen\_card} vs.
          \texttt{lost\_or\_stolen\_phone}, or various
          \texttt{top\_up\_*} variants).
\end{itemize}

\subsection{Classification Report}

\begin{figure}[H]
    \centering
    \includegraphics[width=0.95\textwidth]{figs/03_classification_report.png}
    \caption{Per-class classification report on the test set, showing precision,
    recall, F1-score, and support for each of the 77 intent labels.}
    \label{fig:cls-report}
\end{figure}

\subsection{Example Predictions from the Notebook Demo}

\begin{table}[H]
    \centering
    \caption{Example predictions from the notebook inference demo cell.
    All five examples below were predicted correctly.}
    \label{tab:examples}
    \begin{tabular}{p{6.4cm}ll}
        \toprule
        \textbf{Input Message}                          & \textbf{True Label}              & \textbf{Predicted}              \\
        \midrule
        ``I want to activate my new card''              & activate\_my\_card               & activate\_my\_card              \\
        ``Why was I charged an extra fee?''             & extra\_charge\_on\_statement     & extra\_charge\_on\_statement    \\
        ``How do I change my PIN number?''              & change\_pin                      & change\_pin                     \\
        ``My card payment is still pending''            & pending\_card\_payment           & pending\_card\_payment          \\
        ``I lost my phone''                             & lost\_or\_stolen\_phone          & lost\_or\_stolen\_phone         \\
        ``What is the exchange rate for USD to EUR?''   & exchange\_rate                   & exchange\_rate                  \\
        \bottomrule
    \end{tabular}
\end{table}

\begin{figure}[H]
    \centering
    \includegraphics[width=0.95\textwidth]{figs/04_inference_demos.png}
    \caption{Notebook inference demo cell output: the fine-tuned model predicts
    the correct intent label for each of the six sample messages.}
    \label{fig:inference-demo}
\end{figure}

\subsection{Standalone Command-Line Inference}

The standalone CLI script \texttt{scripts/inference.py} loads the saved
checkpoint and produces a prediction for a single user message:

\begin{lstlisting}[language=bash,caption={CLI inference command demonstrated in the video.}]
python scripts/inference.py \
    --config configs/inference.yaml \
    --message "I need to change my PIN number"
\end{lstlisting}

\begin{figure}[H]
    \centering
    \includegraphics[width=0.95\textwidth]{figs/05_cli_inference.png}
    \caption{Standalone inference script invoked from the command line. The
    script loads the saved checkpoint, the 77 intent labels, and predicts
    \texttt{change\_pin} for the input
    \textit{``I need to change my PIN number.''}}
    \label{fig:cli-inference}
\end{figure}

% ==============================================================================
\section{Inference Implementation}
\label{sec:inference}
% ==============================================================================

\subsection{IntentClassification Class}

The inference module \texttt{scripts/inference.py} implements the required
\texttt{IntentClassification} class with exactly two main methods, matching
the interface specified by the project brief:

\begin{lstlisting}[language=Python, caption={Required inference class interface.}]
class IntentClassification:
    def __init__(self, model_path):
        """
        Load the YAML configuration file referenced by `model_path`,
        then load the tokenizer, the saved model checkpoint, and the
        label mapping.
        """
        ...

    def __call__(self, message):
        """
        Receive a banking customer message string and return the
        predicted intent label.
        """
        ...
        return predicted_label
\end{lstlisting}

\textbf{Key design decisions:}
\begin{itemize}[noitemsep]
    \item \texttt{model\_path} points to a YAML configuration file
          (\texttt{configs/inference.yaml}) that specifies the checkpoint
          directory, the label-map JSON, and generation settings (temperature,
          max new tokens, etc.).
    \item The model is loaded in 4-bit quantization mode, identical to
          training, so memory usage is minimized.
    \item Greedy decoding (\texttt{do\_sample=False}, \texttt{temperature=0})
          is used so that predictions are deterministic.
    \item Fuzzy label matching (\texttt{difflib.get\_close\_matches}) is used
          as a safety net: if the generated text is not exactly one of the 77
          known labels, the closest match (above a similarity cutoff of 0.4)
          is chosen.
    \item The CLI script also supports \texttt{--evaluate} (full test-set
          evaluation) and \texttt{--interactive} (REPL) modes.
\end{itemize}

\textbf{Usage example:}
\begin{lstlisting}[language=Python, caption={Programmatic usage of the IntentClassification class.}]
from scripts.inference import IntentClassification

classifier = IntentClassification("configs/inference.yaml")
label = classifier("I need to change my PIN")
print(label)  # -> "change_pin"
\end{lstlisting}

\subsection{Inference Configuration File}

The YAML configuration referenced by \texttt{model\_path} contains:

\begin{lstlisting}[language={},caption={\texttt{configs/inference.yaml}}]
model:
  name: "unsloth/Llama-3.2-3B-Instruct-bnb-4bit"
  checkpoint_path: "saved_model"
  max_seq_length: 512
  load_in_4bit: true
  dtype: null

data:
  label_map_file: "sample_data/label_map.json"
  test_file: "sample_data/test.csv"

generation:
  max_new_tokens: 64
  temperature: 0.0
  do_sample: false
\end{lstlisting}

% ==============================================================================
\section{Source Code Organization}
\label{sec:source}
% ==============================================================================

All source code is hosted publicly on GitHub at:
\begin{center}
    \repolink{}
\end{center}

The repository follows the structure required by the project brief:

\begin{lstlisting}[language={}, caption={Repository structure.}]
banking-intent-unsloth/
|-- scripts/
|   |-- preprocess_data.py     # download + sample BANKING77
|   |-- train.py               # Unsloth + SFTTrainer fine-tuning
|   |-- inference.py           # IntentClassification class
|-- configs/
|   |-- train.yaml             # training hyperparameters
|   |-- inference.yaml         # inference settings
|-- sample_data/
|   |-- train.csv              # 3,075 sampled training rows
|   |-- test.csv               # 770 sampled test rows
|   |-- label_map.json         # id <-> label mapping
|-- Banking77_Unsloth_Colab.ipynb   # autopilot notebook (with outputs)
|-- train.sh                   # one-click training pipeline
|-- inference.sh               # one-click inference / demo
|-- requirements.txt           # Python dependencies
|-- README.md                  # full setup, train, run, results
|-- Report_Lab02.tex           # this report
|-- .gitignore
\end{lstlisting}

\begin{figure}[H]
    \centering
    \includegraphics[width=0.95\textwidth]{figs/06_github_repo.png}
    \caption{Public GitHub repository, showing the project structure required
    by the assignment, plus the README rendered with the actual results.}
    \label{fig:github-repo}
\end{figure}

\begin{itemize}[noitemsep]
    \item \textbf{scripts/preprocess\_data.py:} downloads BANKING77 from
          Hugging Face (with three-tier fallback for new \texttt{datasets}
          releases), samples a stratified subset, cleans text, and saves CSV
          plus label-map files.
    \item \textbf{scripts/train.py:} loads the model with Unsloth, attaches
          LoRA adapters, trains with \texttt{trl.SFTTrainer} +
          \texttt{SFTConfig}, saves the checkpoint, and evaluates on the test
          set.
    \item \textbf{scripts/inference.py:} the \texttt{IntentClassification}
          class plus a CLI front-end with three modes: single message,
          full evaluation, interactive REPL.
    \item \textbf{configs/}: declarative YAML configurations for all
          hyperparameters and paths.
    \item \textbf{train.sh / inference.sh}: convenience shell scripts.
    \item \textbf{Banking77\_Unsloth\_Colab.ipynb:} the autopilot notebook
          that orchestrates the entire pipeline (data $\rightarrow$ training
          $\rightarrow$ evaluation $\rightarrow$ inference $\rightarrow$
          GitHub push). All outputs are preserved in the committed notebook
          for reproducibility evidence.
\end{itemize}

% ==============================================================================
\section{Video Demonstration}
\label{sec:video}
% ==============================================================================

A short screen-recorded video (approximately 2--3 minutes) demonstrates the
inference behaviour of the trained model. The video covers the four points
required by the assignment:
\begin{enumerate}[noitemsep]
    \item Live execution of the standalone inference script
          (\texttt{python scripts/inference.py ...}).
    \item At least one example input message (\textit{``I need to change my
          PIN number''}).
    \item The intent label predicted by the model
          (\texttt{change\_pin}).
    \item The final accuracy on the test set
          (\textbf{0.8792}, $677 / 770$).
\end{enumerate}

\textbf{Video link (Google Drive, public):}\\
\videolink{}

The link is also embedded in the project README on GitHub, so the grader can
access everything from the public repository alone.

% ==============================================================================
\section{Discussion}
\label{sec:discussion}
% ==============================================================================

\subsection{Why a Generative Approach?}

The project brief asks for a sequence-classification fine-tune. Two natural
implementations exist:
\begin{enumerate}[noitemsep]
    \item Attach a classification head on top of an encoder
          (e.g., DistilBERT, RoBERTa).
    \item Use a causal LLM with SFT and let the model generate the intent
          label as text.
\end{enumerate}

We chose option (2) because the assignment explicitly mandates the
\textbf{Unsloth} workflow, and Unsloth's primary documented pipeline is
SFT on causal LMs. With a strong instruction-following base model
(LLaMA 3.2 3B Instruct), an Alpaca-style prompt converges quickly and the
final accuracy of \textbf{87.92\%} is competitive with classifier-head
baselines on similarly small subsets of BANKING77.

\subsection{Why 87.92\% Instead of Higher?}

Manual inspection of the per-class report (Figure~\ref{fig:cls-report}) shows
that the bulk of errors come from semantically overlapping classes such as:
\begin{itemize}[noitemsep]
    \item \texttt{pending\_card\_payment} vs. \texttt{pending\_transfer}
    \item \texttt{lost\_or\_stolen\_card} vs. \texttt{compromised\_card}
    \item Various \texttt{top\_up\_*} variants
    \item Various \texttt{transfer\_*} and \texttt{cash\_withdrawal\_*}
          variants
\end{itemize}
These are intrinsically hard distinctions for a model trained on only 40
examples per class. With more training data, larger LoRA rank, or longer
training, accuracy is expected to climb above 90\%.

\subsection{Engineering Lessons}

Several practical issues had to be solved while building this pipeline:
\begin{itemize}[noitemsep]
    \item \textbf{Hugging Face \texttt{datasets} v3+ removed support for
          dataset scripts.} The original BANKING77 repo had a Python loader,
          which broke on Colab's latest environment. We added a three-tier
          fallback (default $\rightarrow$ parquet branch $\rightarrow$ direct
          parquet URLs).
    \item \textbf{TRL renamed several \texttt{SFTTrainer} arguments into
          \texttt{SFTConfig}.} \texttt{dataset\_text\_field},
          \texttt{max\_seq\_length}, \texttt{packing}, and
          \texttt{dataset\_num\_proc} are now \texttt{SFTConfig} fields.
          We migrated to the new API.
    \item \textbf{Unsloth and Hugging Face create large cache folders by
          default.} Initially these polluted the GitHub push. We redirected
          all caches to \texttt{/content/hf\_cache} via
          \texttt{HF\_HOME} / \texttt{HUGGINGFACE\_HUB\_CACHE} environment
          variables, and also added the cache directories to
          \texttt{.gitignore}.
    \item \textbf{Re-running the notebook in the same Colab session was
          unsafe} because Step~1 deletes and re-clones the project folder;
          if the working directory was inside that folder, Python would lose
          its cwd. We fixed this by chdir-ing back to \texttt{/content} before
          deleting the folder.
\end{itemize}

\subsection{Future Work}

\begin{itemize}[noitemsep]
    \item Experiment with
          \texttt{FastLanguageModel.for\_sequence\_classification} for a
          dedicated classification head, which should give faster inference.
    \item Train on the full BANKING77 training set (10{,}003 samples) and a
          higher LoRA rank (e.g., $r=32$).
    \item Apply data augmentation (paraphrasing, back-translation) to the
          underrepresented classes.
    \item Compare with smaller encoder baselines (DistilBERT, RoBERTa, MPNet).
\end{itemize}

% ==============================================================================
\section{Conclusion}
\label{sec:conclusion}
% ==============================================================================

We successfully fine-tuned the \textbf{LLaMA 3.2 3B Instruct} model with
\textbf{QLoRA} via \textbf{Unsloth} on a sampled subset of the BANKING77
dataset, reaching \textbf{87.92\%} test accuracy on a 77-class intent
classification task in only \textbf{4.1 minutes} of training on a single H100
GPU. The deliverables -- a public GitHub repository with the required folder
structure, an autopilot Colab notebook with all outputs preserved, a
standalone \texttt{IntentClassification} class with the required
\texttt{\_\_init\_\_}/\texttt{\_\_call\_\_} signature, and a public demo
video -- collectively satisfy every item of the project rubric.

% ==============================================================================
% References
% ==============================================================================
\newpage
\begin{thebibliography}{9}

\bibitem{casanueva2020}
Casanueva, I., Tem\v{c}inas, T., Gerber, D., Fitzpatrick, M., \& Vuli\'c, I. (2020).
\textit{Efficient Intent Detection with Dual Sentence Encoders}.
Proceedings of the 2nd Workshop on NLP for ConvAI -- ACL 2020.
\url{https://arxiv.org/abs/2003.04807}

\bibitem{hu2021lora}
Hu, E.~J., Shen, Y., Wallis, P., Allen-Zhu, Z., Li, Y., Wang, S., \& Chen, W. (2021).
\textit{LoRA: Low-Rank Adaptation of Large Language Models}.
arXiv preprint arXiv:2106.09685.
\url{https://arxiv.org/abs/2106.09685}

\bibitem{unsloth}
Unsloth AI. (2024). \textit{Unsloth -- Fine-tune LLMs 2-5x faster with 80\% less memory}.
\url{https://github.com/unslothai/unsloth}

\bibitem{dettmers2023qlora}
Dettmers, T., Pagnoni, A., Holtzman, A., \& Zettlemoyer, L. (2023).
\textit{QLoRA: Efficient Finetuning of Quantized Language Models}.
arXiv preprint arXiv:2305.14314.
\url{https://arxiv.org/abs/2305.14314}

\bibitem{llama3}
Meta AI. (2024). \textit{Llama 3 Model Card}.
\url{https://github.com/meta-llama/llama3}

\bibitem{trl}
von Werra, L., et al. \textit{TRL -- Transformer Reinforcement Learning}.
\url{https://github.com/huggingface/trl}

\end{thebibliography}

\end{document}
